% Source: https://artofproblemsolving.com/community/c3h1993433_modular_exponentiation_why_all_different
% Stable identity: aops-topic-1993433-q01
(a) Since $\gcd(k,p-1)=1$, there is an integer $u$ with $1\le u\le p-1$ and $ku\equiv1\pmod{p-1}$; write $ku=1+(p-1)v$ with $v\ge0$. For $x\not\equiv0\pmod p$, Fermat's little theorem gives $x^{p-1}\equiv1$, so
\[
(x^k)^u=x^{ku}=x\cdot(x^{p-1})^v\equiv x\pmod p .
\]
Hence if $x^k\equiv y^k$ for nonzero residues $x,y$, raising to the power $u$ gives $x\equiv y$. So $x\mapsto x^k$ is one-to-one on the $p-1$ nonzero residues (it never gives $0$, since $p\nmid x$), and a one-to-one map of a finite set to itself is a permutation. $\blacksquare$

(b) For $p=5$, $k=3$: $3\cdot3=9=1+4\cdot2$, so cubing twice gives $x^9\equiv x$; cubing undoes itself, so it must be a permutation, and indeed $1\mapsto1$, $2\mapsto3$, $3\mapsto2$, $4\mapsto4$. Squaring fails because $\gcd(2,4)\ne1$: $1^2\equiv4^2\equiv1$ and $2^2\equiv3^2\equiv4$, so only $1$ and $4$ are hit.
